%============================================================ 7. 고찰 및 결론 (압축판)
\section{고찰 및 결론}
\label{sec:discussion}

\subsection{결과의 해석}

두 계층의 기여가 가산적이라는 것(\ref{sec:res_g3}절)은 지휘관 모델 교체와 기동 정책 재학습을
서로 기다리지 않고 진행할 수 있다는 운용상의 자유를 뜻한다. 지휘관의 성능을 설명한 것이
커버리지가 아니라 계획 안정성이라는 관찰(\ref{sec:res_mechanism}절)은 주기적으로 재계획하는 계층
구조 일반에 해당한다. 상위 결정의 변경은 하위 계층에 회수되지 않는 이동과 그물을 남기므로, 매
호출이 독립적인 언어모델을 계획 계층에 둘 때는 계획 품질을 시간에 걸친 일관성으로 평가하고 상태
표현이 그 일관성을 지원해야 한다. 가장 큰 개선이 총 선회량에서 나타난 것은 경유점이 격자에
고정되고 유효 마스크가 무의미한 위치를 배제하기 때문이며, 유효 셀 안에서 점수맵은 거의 단봉이므로
이득은 분포의 다봉성이 아니라 마스크와 격자 양자화에서 온다.

\subsection{한계와 향후 과제}
\label{sec:limitations}

검증은 시뮬레이션에 한정되며 그물 거동·얽힘 판정·관측이 단순화되어 있다. 관측 불확실성은 배정이
딛고 선 무리 구성을 흔들므로 지휘관 계층에 더 크게 작용할 것이다. 규모가 적 10척 대 방어정 3척으로
고정이고 적이 규칙 기반이라, 배정 조합이 급증하는 규모와 그물 배치를 학습해 회피하는 적응적
적은 범위 밖이며, 연속 행동공간 기준선과 구성요소별 절제도 과제로 남는다.

\subsection{결론}
\label{sec:conclusion}

적보다 느린 방어정으로 자폭 군집을 막는 2계층 체계를 제안하고, 시드 대응 $2\times2$ 설계의
720 에피소드로 두 계층의 기여를 분리 측정하였다. 제안 체계는 같은 방어 수준에서 자원과 기동
비용을 크게 줄였으며, 두 계층의 기여는 가산적이었고, 언어모델 지휘관의 이득은 능력 임계 위에서만
나타났다. 실해역 이식과 관측 불확실성 하에서의 계획 안정성은 향후 과제로 남는다.
